\documentclass[../main.tex]{subfiles}

\begin{document}

\section{Main concepts}

In this section, we present the generalized maps, graph transformation rules and their subsequent concepts which are necessary to the understanding of our contribution.

\subsection{Generalized maps}

\emph{Generalized maps} (or G-maps) \cite{lienhardt1991topological,
  Damiand-Lienhardt14} allow the representation of manifold geometric objects
(with or without boundaries), based on some cellular $n$-dimensional topological
structure.

\begin{figure*}[t]
    \begin{subfigure}{.2\textwidth}
        \centering
        \includegraphics[height=.2\textheight]{house-ABCDE.pdf}
        \caption{}
        \label{fig:house-ABCDE}
    \end{subfigure}
    \hfill
    \begin{subfigure}{.2\textwidth}
        \centering
        \includegraphics[height=.2\textheight]{house-faces-decomposition}
        \caption{}
        \label{fig:house-faces-decomposition}
        \end{subfigure}
    \hfill
    \begin{subfigure}{.2\textwidth}
        \centering
        \includegraphics[height=.2\textheight]{house-show-1-links.pdf}
        \caption{}
        \label{fig:house-show-1-links}
    \end{subfigure}
    \hfill
    \begin{subfigure}{.2\textwidth}
        \centering
        \includegraphics[height=.2\textheight]{house-show-all-links.pdf}
        \caption{}
        \label{fig:house-show-all-links}
    \end{subfigure}
    \caption{Cell decomposition of a geometric 2D object}
    \label{fig:house-cell-decomposition}
\end{figure*}

The representation of an object as a G-map comes intuitively from its
decomposition into topological cells (vertices, edges, faces, volumes, and so
on). For example, the 2D topological object (Fig.~\ref{fig:house-ABCDE}) can be
decomposed into a $2$-dimensional G-map (or $2$-G-map). The object is first
decomposed into faces (Fig.~\ref{fig:house-faces-decomposition}). These faces
are \emph{linked} along their common edge with an $\alpha 2$ link, drawn in
blue. Each edge on the border of the object is connected to itself (the blue
$\alpha 2$ link forms a loop). The index "$2$" means that the link connect two
$2$-dimensional (possibly a single one) faces. In order to alleviate the
reading, $\alpha 2$ links will not be drawn in every subsequent figure. In the
same way, faces are split into edges connected with red $\alpha 1$ links
(Fig~\ref{fig:house-show-1-links}). At last, edges themselves are split into
vertices with black $\alpha 0$ links (Fig~\ref{fig:house-show-all-links}), to
produce the $2$-G-map describing the objects shown in
Fig.~\ref{fig:house-ABCDE}. A G-map is therefore a graph, the nodes of which are
the vertices (named \emph{darts}) obtained at the end of the process and the
edges are $\alpha i$ links.

G-maps have conditions guaranteeing objects consistency, for example, two faces
are always linked along an edge.

\begin{figure*}[t]
  \begin{subfigure}{.19\textwidth}
    \centering
    \includegraphics[height=.2\textheight]{figures/house-2-cell-a.pdf}
    \caption{}
    \label{fig:house-2-cell-a}
  \end{subfigure}
  \hfill
  \begin{subfigure}{.19\textwidth}
    \centering
    \includegraphics[height=.2\textheight]{figures/house-1-cell-a.pdf}
    \caption{}
    \label{fig:house-1-cell-a}
  \end{subfigure}
  \hfill
  \begin{subfigure}{.19\textwidth}
    \centering
    \includegraphics[height=.2\textheight]{figures/house-1-cell-c.pdf}
    \caption{}
    \label{fig:house-1-cell-c}
  \end{subfigure}
  \hfill
  \begin{subfigure}{.19\textwidth}
    \centering
    \includegraphics[height=.2\textheight]{figures/house-0-cell-a.pdf}
    \caption{}
    \label{fig:house-0-cell-a}
  \end{subfigure}
  \hfill
  \begin{subfigure}{.19\textwidth}
    \centering
    \includegraphics[height=.2\textheight]{figures/house-1-cell-a0-a.pdf}
    \caption{}
    \label{fig:house-1-cell-a0-a}
  \end{subfigure}

  \caption{Orbit decomposition of a geometric 2D object}
  \label{fig:house-orbit-decomposition}
\end{figure*}

Topological cells are not explicitly represented in G-maps but only implicitly
defined as subgraphs. They can be computed using graph traversals defined by an
originating node and by a given set of link labels named \emph{orbit}. For
example, the 2-cell adjacent to some dart $a$ (or the object face matching $a$)
(Fig.~\ref{fig:house-2-cell-a}) is the subgraph which contains $a$ and all darts
reachable from $a$, using links labeled by $\alpha 0$ or $\alpha 1$
(\textit{i.e.} darts $a$, $b$, $c$, $d$, $e$, $f$, $g$ and $h$) and the links
themselves. This subgraph is denoted by $G\orbit{\alpha 0 \alpha 1}(a)$ and
models the face $BCED$ (Fig.~\ref{fig:house-ABCDE}).
$\orbit{\alpha 0, \alpha 1}$ is the \emph{type} of the orbit. The $1$-cell
adjacent to $a$ (or the object edge matching to $a$)
(Fig.~\ref{fig:house-1-cell-a}) is the subgraph $G \orbit{\alpha 0 \alpha 2}(a)$
which contains $a$ and all reachable darts using links labeled by $\alpha 0$ or
$\alpha 2$ (\textit{i.e.} darts $a, b, l$ and $m$) and the corresponding links.
It represents the topological edge $BC$. Thanks to $\alpha 2$-loops on the
object border, $G\orbit{\alpha 0, \alpha 2}(c)$ (Fig.~\ref{fig:house-1-cell-c}
also represents the edge $CE$ matching dart $c$. The $0$-cell adjacent to $a$
(or the object vertex matching dart $a$) (Fig.~\ref{fig:house-0-cell-a}) is the
subgraph $G\orbit{\alpha 1, \alpha 2}(a)$ and represents vertex $B$. Note that
orbits are more general than cells. For example, the face edge
$G\orbit{\alpha 0}(a)$ (Fig.~\ref{fig:house-1-cell-a0-a}) is an
$\orbit{\alpha 0}$ orbit incident to $a$.


\subsection{Graph transformation rules}

Jerboa \cite{belhaouari2014jerboa} graph transformation rules allow to
formalize operations over G-maps.

\begin{figure}[t]
  \centering
  \maxsizebox{\columnwidth}{.15\columnwidth}{\includegraphics{triangulation}}
  \caption{Rule $L \longrightarrow R$ of face triangulation}
  \label{fig:rule-triangulation}
\end{figure}

\begin{figure}[t]
  \centering
  \begin{subfigure}[t]{.45\columnwidth}
    \centering
    \includegraphics[height=.25\textheight]{house-triangulated-left-match}
    \caption{Graph $G$}
    \label{fig:house-triangulation-left-match}
  \end{subfigure}
  \hfill
  \begin{subfigure}[t]{.45\columnwidth}
    \centering
    \includegraphics[height=.25\textheight]{house-triangulated-right-match}
    \caption{Graph $H$}
    \label{fig:house-triangulation-right-match}
  \end{subfigure}
  \caption{Application of the triangulation rule of figure \ref{fig:rule-triangulation} on matched dart $a_{0}$}
  \label{fig:before-after-house-triangulation}
\end{figure}

In a few words, a \emph{rule} $r : L \longrightarrow R$ and a \emph{match}
$m : L \to G$ to a G-map $G$, describe the transformation
$G \longrightarrow^{r,m} H$ from $G$ to $H$. $m$ allows to replace a sub-graph
of $G$ described by the left-hand side of the rule $L$ with another one
described the right-hand side $L$, on order to produce $H$.

\maxime{Introduire terme d'orbite dans une règle}

Informally, in the triangulation rule shown in
Fig.~\ref{fig:rule-triangulation}, the left side is made of only one node
$n_{0}$ labeled with the face $\orbit{\alpha 0, \alpha 1}$ type: This way, it
allows to match any face. For the match $m : n0 \mapsto a_{0}$ from $L$ to $G$
(Fig.~\ref{fig:house-triangulation-left-match}), the whole face
$G\orbit{\alpha 0, \alpha 1}(a_{0})$ is matched, highlighted in orange in the
figure. $alpha 0$ (resp. $\alpha 1$) links are represented by black (resp. blue)
segments. On the right side, the node $n_{0}$ has label $\orbit{\alpha 0, \_}$.
This means that when applying this rule, $\alpha 0$ links of nodes $n_{0}$ have
been preserved, while $\alpha 1$ links have been deleted. Hence, in $H$
(Fig.~\ref{fig:house-triangulation-right-match}), the $\alpha 0$-links of the
matched orange face have been preserved while $\alpha 1$-links have been
deleted. In other words, the edges of the orange face are disconnected in $H$.
The new node $n_{2}$ in $H$ creates darts copied from the matched face. This is
why there are eight blue darts ($a_2$, $b_2$ \dots $h_2$) created from orange
darts ($a_0$, $b_0$ \dots $h_0$). Because $n_{2}$ is labeled
$\orbit{\alpha 1, \alpha 2}$, the orange $\alpha 0$ (resp. $\alpha 1$) links on
the left side of the rule are relabeled to $\alpha 1$ (resp. $\alpha 2$) links
in the blue copy. Therefore, the rule creates a dual vertex to the matched face.
Finally, the node $n_{1}$ creates a green copy with eight darts ($a_1$, $b_1$
\dots $h_1$), deletes the left-side $\alpha 0$ links and relabels the left-side
$\alpha 1$ links to $\alpha 2$ links.

All the highlighted links (Fig.~\ref{fig:before-after-house-triangulation}) are
referred as \emph{implicit links} in the rule nodes. Conversely, the links
connecting the nodes together in $R$ are referred as \emph{explicit links}. For
example, the explicit $\alpha 0$ links between $n1$ and $n2$ link one-to-one the
green and blue darts. Therefore, the rule creates the four new edges
$\orbit{\alpha 0, \alpha 2}$ in the triangulation of the matched square face. In
the same way, the explicit $\alpha 1$ links between $n0$ and $n1$ link
one-to-one the orange and green darts.

The $n0$ node (Fig.~\ref{fig:rule-triangulation}) is a \emph{preserved node}
because it belongs to both the left and right side of the rule. Nodes $n1$ and
$n2$ are \emph{created nodes} because they belong only to the right side of the
rule. This rule does not comprise any \emph{deleted node} because none of its
nodes belong only to the left side.

Additionally, in the triangulation rule, the $\orbit{\alpha 0, \alpha 1}$-orbit
incident to the node $n_{0}$ is said to be complete since each node in the orbit
matches either explicitly or implicitly both an $\alpha 0$-link and an
$\alpha 1$-link. As a result, the matched orbit $G\orbit{o}(a_{0})$ is a face,
i.e. every darts are incident to both an $\alpha 0$-link and an $\alpha 1$-link.
Conversely, the $\orbit{\alpha 0, \alpha 2}$-orbit incident to $n_{0}$ is
incomplete since the $\alpha 2$-link is neither explicitly nor implicitly
matched by $n_{0}$. As a result, the matched orbit $\orbit{\alpha 0}$ is a face
edge and, therefore, is different from $G\orbit{\alpha 0, \alpha 2}(a_{0})$
since $\alpha 2$-links are not matched.

Jerboa's rules provide with syntactic properties guaranteeing G-maps consistency
preservation.

%%% Local Variables:
%%% mode: latex
%%% TeX-master: "../main"
%%% End:
